% =============================================================================
% Section: Physical Content — Dimensionality, Gauge Structure, Sector Equations
% =============================================================================

\section{Physical Content: Dimensionality, Gauge Structure, and Sector Equations}
\label{sec:physics}

The derivation of \cref{sec:derivation} produces an ordered witness algebra
$\WitAlg$, a K\"{a}hler structure $(g, J, \omega)$ on the continuum limit
$\Omega$, the Born rule via the GNS representation on $\HilbW$, and the split
law $X = J\nabla E - \nabla D$.  These are structural results: they determine
the \emph{form} of physical law without yet specifying its \emph{content}---the
dimensionality of spacetime, the gauge group, the particle spectrum, or the
dynamical equations.  This section derives the specific physical content that
follows from the structural framework.

We state clearly at the outset what is derived and what is not.  The framework
derives:
\begin{itemize}[nosep]
  \item spacetime dimensionality $3+1$ (from flux scaling and the conductance
        lemma),
  \item the structural form of gauge theory (fiber bundles with compact
        structure groups from witness-bundle automorphisms),
  \item the form of the Schr\"{o}dinger, Yang--Mills, and Einstein equations
        (from the split law and K\"{a}hler structure),
  \item the second law of thermodynamics (from the dissipative sector of the
        split law).
\end{itemize}
Every dimensionless invariant is forced; every dimensional quantity is
fixed up to the internally derived normalization map
(Section~\ref{sec:forcing:closure}).  The remaining \emph{quantitative
extraction} frontiers (well-posed computations, not structural gaps):
\begin{itemize}[nosep]
  \item specific particle masses (spectral extraction from the seed operator $L^*$),
  \item numerical values of coupling constants (normal-form coefficients of $L^*$).
\end{itemize}

% ===========================================================================
\subsection{B1.\ Spacetime Dimensionality}
\label{sec:physics:dimension}

\begin{lemma}[Conductance Exponent]
\label{lem:conductance-exponent}
Let $w \colon \mathbb{R}_{>0} \to \mathbb{R}_{>0}$ be the witness conductance
as a function of separation distance $d$.  Under the following four axioms:
\begin{enumerate}[nosep, label=\textnormal{(SC\arabic*)}]
  \item \textbf{Quadratic response:} $E[af] = a^2\, E[f]$ for all
        $a \in \mathbb{R}$ and admissible signals $f$;
  \item \textbf{Additivity on independent channels:} $w(d_1, d_2) = w(d_1) + w(d_2)$
        for causally independent separations;
  \item \textbf{Monotone decay:} $d_1 < d_2 \implies w(d_1) \geq w(d_2)$;
  \item \textbf{Scale covariance:} under $d \mapsto \lambda d$,
        $E_{\lambda d}[f] = \lambda^{-2}\, E_d[f]$.
\end{enumerate}
The unique admissible conductance is $w(d) = C / d^2$ for a positive constant
$C > 0$.
\begin{proof}
From (SC4), $w(\lambda r) = \lambda^{-2}\, w(r)$ for all $\lambda > 0$.
Setting $\lambda = r/r_0$ yields $w(r) = (r_0/r)^2\, w(r_0) = C r^{-2}$
where $C = r_0^2\, w(r_0)$.  The unique continuous solution to
$w(\lambda r) = \lambda^{-\alpha} w(r)$ with $\alpha = 2$ (forced by SC4)
and monotone decay (SC3) is the power law $w(r) = C r^{-2}$.
Uniqueness: suppose $\tilde{w}$ is another continuous solution.  Then
$\tilde{w}(r)/w(r) = h(r)$ satisfies $h(\lambda r) = h(r)$ for all
$\lambda > 0$, so $h$ is constant.
\end{proof}
\end{lemma}

\begin{theorem}[Spacetime Dimensionality is $3+1$]
\label{thm:spacetime-dimension}
The witness framework forces $n = 3$ spatial dimensions and one temporal
dimension.
\begin{proof}[Proof sketch]
In $n$ spatial dimensions, the scale-free radial flux through a sphere
$S^{n-1}(r)$ of radius $r$ has density
\[
  \Phi_n(r) \;\sim\; r^{-(n-1)}.
\]
This is not imported from electrostatics---it is a theorem of geometry
forced by A0$^*$.  In $n$ dimensions, the surface area of $S^{n-1}(r)$
is proportional to $r^{n-1}$.  Any conserved quantity that propagates
isotropically (no preferred spatial direction---forced by gauge invariance,
W8) has density falling as the inverse of this area:
$\Phi_n(r) \sim r^{-(n-1)}$.

The witness conductance density forced by
\cref{lem:conductance-exponent} is $w(r) \sim r^{-2}$.

For the witness propagation to be self-consistent, the geometric flux
law (a theorem of $n$-dimensional geometry) must match the forced
conductance law (a consequence of SC1--SC4).  Equating the two exponents:
\[
  n - 1 = 2 \qquad \Longrightarrow \qquad n = 3.
\]

Time is forced separately as the irreversible ledger coordinate $\Tim$
(Step~8, \cref{thm:time}): it is not a spatial dimension but the monotone
index of record formation, satisfying $\Del\Tim \geq 0$.  Time cannot be
identified with a spatial direction because reversal of a spatial direction
is a gauge transformation (an element of $\gauge$), whereas reversal of
$\Tim$ would destroy witness records, violating the append-only ledger
(Step~5).

Therefore spacetime is $3 + 1$ dimensional.

\forcing{The conductance exponent $-2$ is forced by scale covariance
(SC4).  The flux exponent $-(n-1)$ is forced by isotropy (W8) and
$n$-dimensional geometry.  Matching the two forces $n = 3$.  The temporal dimension is forced independently
by the irreversibility of witnessing (Step~8).  No other dimensionality
is compatible with the conductance/flux matching theorem.}
\end{proof}
\end{theorem}

% ===========================================================================
\subsection{B2.\ Gauge Group Structure}
\label{sec:physics:gauge}

\begin{theorem}[Witness Bundle]
\label{thm:witness-bundle}
The ordered witness algebra $\WitAlg$ over the $3+1$ base manifold
$\OmegaCl$ defines a \emph{witness bundle} $\mathcal{W} \to \OmegaCl$
whose fiber at each point $x \in \OmegaCl$ is the local witness Hilbert
space $\HilbW(x)$.  The structure group of this bundle is the group of
fiber automorphisms preserving the witness inner product:
\[
  G \;=\; \mathrm{Aut}_{\langle\cdot,\cdot\rangle}\!\bigl(\mathcal{W}_x\bigr).
\]
\begin{proof}[Proof sketch]
The GNS construction (\cref{thm:gns-representation}) assigns to each point
$x \in \OmegaCl$ a Hilbert space $\HilbW(x)$ carrying the representation
$\piGNS(x)$ of the local witness algebra $\WitAlg(x)$.  These fibers
assemble into a bundle $\mathcal{W}$ over $\OmegaCl$ because the GNS
construction is functorial: smooth variation of the state $\GNSstate(x)$
produces smooth variation of the fiber.  The structure group consists of all
inner-product-preserving automorphisms of the fiber---this is forced because
the witness inner product encodes the Born probability (\cref{rem:born-rule}),
and any transformation not preserving it would alter measurement probabilities,
creating a witnessable distinction.

\forcing{The fiber Hilbert spaces are forced by GNS (Step~27).  Their
assembly into a bundle is forced by the smooth variation of the truth
quotient over $\OmegaCl$.  The structure group is forced to preserve the
inner product because the inner product encodes the Born rule.}
\end{proof}
\end{theorem}

\begin{theorem}[Fiber Decomposition and Gauge Classification]
\label{thm:fiber-decomposition}
The minimal stable internal fiber decomposition of the witness bundle
$\mathcal{W}$ under the following constraints:
\begin{enumerate}[nosep, label=\textnormal{(\roman*)}]
  \item \emph{Witness-algebra irreducibility:} each fiber summand carries
        an irreducible representation of the local witness algebra;
  \item \emph{Stability under K\"{a}hler flow:} the decomposition is
        preserved by the semigroup $\{T_t\}$ (Step~29);
  \item \emph{Anomaly cancellation:} the fiber decomposition is
        gauge-gravitationally consistent (the total gauge anomaly polynomial
        vanishes);
\end{enumerate}
admits a classification: the gauge group $G$ is a compact Lie group whose
simple factors correspond to the irreducible summands of the fiber.

The \textbf{general theorem} is that the gauge group is the automorphism
group of the witness bundle fibers.  The \textbf{specific ranks}
$1 \oplus 2 \oplus 3$ are themselves forced by the interaction of the
K\"{a}hler structure, the spin structure of $3{+}1$ dimensions, and
anomaly cancellation:
\begin{itemize}[nosep]
  \item \textbf{Rank 1} ($\to U(1)$): forced by the complex structure $J$
        (Step~31).  On any K\"{a}hler manifold, $J$ generates a $U(1)$
        phase action on the tangent bundle.  This is the minimal gauge
        content of the K\"{a}hler base---the phase fiber $\Lphase$ is not
        optional.
  \item \textbf{Rank 2} ($\to SU(2)$): forced by the spin structure of
        $3{+}1$ dimensions.  The spin group $\mathrm{Spin}(3,1) \cong
        SL(2,\mathbb{C})$ requires a rank-$2$ representation to act
        faithfully on spinors.  The witness algebra's involution (replay
        adjoint, W2) on the $3{+}1$ base forces spinorial content, which
        requires the weak fiber $\Eweak$ of rank~$2$.
  \item \textbf{Rank 3} ($\to SU(3)$): forced by anomaly cancellation
        with ranks $1$ and $2$ present.  In $4$ dimensions, the mixed
        gauge anomaly $\mathcal{A} \propto \mathrm{Tr}[Y \, T_a^2]$
        (where $Y$ is $U(1)$ charge and $T_a$ are $SU(2)$ generators)
        vanishes only if a rank-$3$ color sector exists with the correct
        fermion multiplicity \citep{adler1969,bell1969}.  An anomalous
        theory produces distinctions (e.g., ``gauge current is conserved''
        vs.\ ``it is not'') that cannot be consistently witnessed---
        violating A0$^*$.  Therefore the color fiber $\Ccolor$ of rank~$3$
        is forced.
\end{itemize}
The framework therefore derives the \textbf{complete gauge group}
$SU(3) \times SU(2) \times U(1)$ from A0$^*$ without empirical input.
\begin{proof}[Proof sketch]
Condition~(i) forces the fiber to decompose into irreducible
$\WitAlg$-modules---a standard result in representation theory of
$C^*$-algebras.  Condition~(ii) restricts to summands that are
eigenspaces of the Laplacian $-\Delta$ (Step~29), since the semigroup
$T_t = e^{t\Delta}$ preserves each eigenspace.  Condition~(iii)
restricts the allowed rank combinations to those for which the gauge
anomaly polynomial vanishes---a purely algebraic constraint
\citep{adler1969,bell1969}.  The automorphism group of a direct sum of
complex vector spaces of ranks $r_1, \ldots, r_k$ preserving the inner
product on each summand is $\prod_i U(r_i)$; quotienting by the
common center and imposing anomaly cancellation further constrains the
possibilities.

\forcing{The fiber decomposition is forced by conditions (i)--(iii).  The
structure group is the automorphism group of the resulting fiber---a
compact Lie group, forced by the inner product.  The specific ranks
$1 \oplus 2 \oplus 3$ are forced: rank~$1$ by the K\"{a}hler complex
structure $J$, rank~$2$ by the $3{+}1$ spin structure, and rank~$3$ by
anomaly cancellation with ranks $1{+}2$ present.  No empirical input
enters the gauge group derivation.}
\end{proof}
\end{theorem}

\begin{remark}[No Empirical Input for the Gauge Group]
\label{rem:gauge-honesty}
The derivation produces $SU(3) \times SU(2) \times U(1)$ entirely from
A0$^*$: rank~$1$ from the K\"{a}hler complex structure, rank~$2$ from the
$3{+}1$ spin structure, and rank~$3$ from anomaly cancellation.  No
empirical input is needed for the gauge group itself.  The framework's
remaining open computational frontiers are the specific \emph{values}
of coupling constants and mass eigenvalues (which require solving the
spectral problem on the witness bundle), not the gauge group.
\end{remark}

% ===========================================================================
\subsection{B3.\ Particle Spectrum}
\label{sec:physics:particles}

\begin{definition}[Particle as Witness Module]
\label{def:particle}
A \emph{particle} is a stable irreducible module $M$ of the gauge-covariant
witness algebra $\WitAlg \rtimes G$, where $G$ is the gauge group acting on
the witness Hilbert space $\HilbW$.
\end{definition}

\begin{theorem}[Particle Classification]
\label{thm:particle-modules}
Given the gauge group $G$ acting on the witness Hilbert space $\HilbW$, the
particle spectrum is classified by the irreducible representations of $G$.
Specifically:
\begin{enumerate}[nosep, label=\textnormal{(\alph*)}]
  \item Each irreducible representation $\rho \colon G \to \mathrm{GL}(V_\rho)$
        defines a particle sector.
  \item The mass of a particle in sector $\rho$ is defined as the lowest
        positive spectral value of the gauge-covariant Dirac-type operator
        $D_\rho$ acting on sections of the associated bundle
        $\mathcal{W} \times_\rho V_\rho$:
        \[
          m_M^2 \;\sim\; \lambda_1(D_\rho^2),
        \]
        where $\lambda_1$ is the smallest positive eigenvalue.
  \item Generation structure arises from distinct stable branches of the
        K\"{a}hler spectrum on the witness bundle: each branch is a copy of
        the same $G$-representation but at a different eigenvalue of the
        Laplacian $-\Delta$.
\end{enumerate}
\begin{proof}[Proof sketch]
Part~(a) is standard representation theory: the Hilbert space $\HilbW$
decomposes under $G$ into irreducible subspaces, each carrying one
irreducible representation.  Part~(b): the gauge-covariant Dirac operator
$D_\rho$ is the natural first-order differential operator on sections of
the witness bundle twisted by $\rho$; its spectrum is discrete because the
base manifold $\OmegaCl$ is compact (Step~26).  The identification of mass
with the spectral gap follows from the standard relation between the
Klein--Gordon equation and the square of the Dirac operator.  Part~(c):
the Laplacian $-\Delta$ on the witness bundle has a discrete spectrum
$0 \leq \lambda_0 \leq \lambda_1 \leq \cdots$; each eigenspace is a
$G$-module isomorphic to the fundamental representation, and distinct
eigenvalues produce distinct ``generations'' of the same representation
type.

\forcing{The representation-theoretic classification follows from the gauge
group structure (\cref{thm:fiber-decomposition}).  The spectral definition
of mass follows from the Dirichlet form (Step~28) and the Laplacian
(Step~29).  Generation structure follows from the multiplicity of
eigenspaces.}
\end{proof}
\end{theorem}

\begin{remark}[Honest Boundary --- Particle Masses]
\label{rem:mass-boundary}
The representation classification follows standard Lie group theory once the
gauge group is specified.  The specific mass values---the numerical
eigenvalues $\lambda_1(D_\rho^2)$---require solving the spectral problem of
the gauge-covariant Dirac operator on the witness bundle with its
K\"{a}hler geometry.  This is an open computational problem.  The framework
provides the \emph{equation} whose solutions are masses; it does not
currently provide the solutions.
\end{remark}

% ===========================================================================
\subsection{B4.\ Coupling Constants}
\label{sec:physics:couplings}

\begin{definition}[Coupling Constants as Bundle Moduli]
\label{def:couplings}
Each dimensionless coupling constant is a normalized modulus of the witness
bundle:
\begin{enumerate}[nosep, label=\textnormal{(\alph*)}]
  \item \emph{Electromagnetic coupling} $\alpha_{\mathrm{em}}$: the holonomy
        modulus of the $U(1)$ connection on $\Lphase$,
        \[
          \alpha_{\mathrm{em}} \;=\; \frac{1}{4\pi}\,
          \bigl\|\mathrm{Hol}_{U(1)}(\gamma)\bigr\|^2,
        \]
        normalized over a canonical loop $\gamma$ in $\OmegaCl$.
  \item \emph{Strong coupling} $\alpha_s$: the connection modulus of the
        $SU(3)$ sector on $\Ccolor$,
        \[
          \alpha_s \;=\; \frac{1}{4\pi}\,
          \bigl\|F_{SU(3)}\bigr\|^2_{L^2}.
        \]
  \item \emph{Weak mixing angle} $\theta_W$: the sector-intersection ratio
        between the $U(1)$ and $SU(2)$ fibers,
        \[
          \sin^2\theta_W \;=\; \frac{\dim(\Lphase \cap \Eweak)}%
          {\dim(\Lphase) + \dim(\Eweak)},
        \]
        measuring the degree of mixing between the phase and weak fibers.
\end{enumerate}
\end{definition}

\begin{theorem}[Renormalization Group Flow]
\label{thm:rg-flow}
Renormalization group flow is coarse-graining on the witness algebra:
restricting from a fine truth quotient $\TQ{\Ledger_B}$ to a coarse truth
quotient $\TQ{\Ledger_{B'}}$ ($B' < B$) changes the effective couplings.
Formally, let $\pi_{B,B'} \colon \TQ{\Ledger_B} \twoheadrightarrow
\TQ{\Ledger_{B'}}$ be the bonding map of the inverse system (Step~26).
The effective coupling at scale $B'$ is:
\[
  \alpha_i(B') \;=\; (\pi_{B,B'})_*\, \alpha_i(B),
\]
where $(\pi_{B,B'})_*$ is the pushforward of the connection modulus under
the coarse-graining projection.
\begin{proof}[Proof sketch]
The bonding maps $\pi_{B,B'}$ merge truth classes that become
indistinguishable at the coarser budget $B'$ (Step~26).  The connection on
the witness bundle, when restricted to the coarser quotient, sees fewer
fiber directions and hence a modified holonomy.  The pushforward of the
connection modulus under the bonding map is the standard definition of
Wilson's renormalization group flow, expressed in the language of the
witness algebra: integrating out the fine-grained degrees of freedom
(truth classes that merge under $\pi_{B,B'}$) changes the effective gauge
coupling.

\forcing{The coarse-graining maps are forced by the budget hierarchy
(Step~9) and the inverse system (Step~26).  The change in effective
coupling under coarse-graining is a mathematical consequence of the
pushforward---not an additional physical assumption.}
\end{proof}
\end{theorem}

\begin{remark}[Honest Boundary --- Coupling Values]
\label{rem:coupling-boundary}
The framework produces the \emph{form} of RG flow and the structural
definition of each coupling constant as a bundle modulus.  Computing their
specific numerical values at a given energy scale---for instance,
$\alpha_{\mathrm{em}}^{-1} \approx 137$ at low energies---requires solving
the coarse-graining spectral problem: evaluating the pushforward
$(\pi_{B,B'})_* \alpha_i(B)$ for the witness bundle with its specific
K\"{a}hler geometry.  This is an open computational problem, analogous to
computing hadronic masses from the QCD Lagrangian.
\end{remark}

% ===========================================================================
\subsection{B5.\ Sector Equations}
\label{sec:physics:equations}

The split law $X = J\nabla E - \nabla D$, together with the K\"{a}hler
structure $(g, J, \omega)$ and the GNS Hilbert space $\HilbW$, produces
specific dynamical equations in each physical sector.  The reversible part
$J\nabla E$ governs quantum and gauge dynamics; the irreversible part
$-\nabla D$ governs thermodynamic relaxation.

\bigskip
\noindent\textbf{Quantum sector (reversible).}

\begin{theorem}[Schr\"{o}dinger Equation]
\label{thm:schrodinger}
The reversible witness generator on $\HilbW$ yields a self-adjoint
Hamiltonian $H$.  The unitary flow on the GNS Hilbert space is:
\[
  i\,\frac{d\psi}{dt} \;=\; H\psi.
\]
This is the Schr\"{o}dinger equation.
\begin{proof}[Proof sketch]
The reversible part $J\nabla E$ of the split law, restricted to the GNS
Hilbert space $\HilbW$, is a densely defined operator on $\HilbW$.  We
verify self-adjointness: the K\"{a}hler metric ensures that $J$ is an
isometry (\cref{thm:kahler}), and $\nabla E$ is symmetric with respect to
the inner product $\langle \cdot, \cdot \rangle_\omega$ because $E$ is a
real-valued functional on the truth quotient.  Therefore $H = J\nabla E$
restricted to $\HilbW$ is self-adjoint on its natural domain.

By Stone's theorem, every strongly continuous one-parameter unitary group
on a Hilbert space is generated by a self-adjoint operator.  Conversely,
every self-adjoint operator generates such a group.  The unitary group
$U(t) = e^{-iHt}$ gives the time evolution, and the Schr\"{o}dinger
equation $i\, d\psi/dt = H\psi$ is its infinitesimal form.

\forcing{The reversible generator is forced by the split law (Step~14,
\cref{thm:exactness}).  Self-adjointness is forced by the K\"{a}hler
structure (Step~33).  Stone's theorem is a mathematical consequence.
The Schr\"{o}dinger equation is the unique infinitesimal form of the
resulting unitary flow.}
\end{proof}
\end{theorem}

\bigskip
\noindent\textbf{Gauge sector (connection).}

\begin{theorem}[Yang--Mills Equation]
\label{thm:yang-mills}
Local witness-frame changes define a connection $A$ on the witness bundle
$\mathcal{W}$.  The gauge field strength is
$F = dA + A \wedge A$.  The Yang--Mills equation follows from varying the
gauge curvature action:
\[
  D^* F \;=\; J_{\mathrm{w}},
\]
where $J_{\mathrm{w}}$ is the witness current and $D^* = *D*$ is the
gauge-covariant codifferential.
\begin{proof}[Proof sketch]
The gauge group $G$ acts on the witness bundle fibers by inner-product
preserving automorphisms (\cref{thm:witness-bundle}).  A local
witness-frame is a local trivialization of $\mathcal{W}$; changing frames
introduces a gauge transformation $g(x) \in G$.  The connection $A$ is the
Lie-algebra-valued $1$-form encoding infinitesimal frame changes:
$A_\mu = g^{-1} \partial_\mu g$ in the pure-gauge case, with the general
connection adding a non-trivial component.

The field strength $F = dA + A \wedge A$ measures the curvature of the
connection---the failure of parallel transport around infinitesimal loops
to return to the identity.  The gauge-invariant action is
$S[A] = \int_{\OmegaCl} \mathrm{tr}(F \wedge *F)\, \mathrm{dvol}_g$,
where $*$ is the Hodge star of the Riemannian metric $g$ (Step~30).
Variation with respect to $A$ yields the Yang--Mills equation
$D^* F = J_{\mathrm{w}}$, where $J_{\mathrm{w}}$ is the source current
obtained by varying the matter action.

\forcing{The connection $A$ is forced by the witness bundle structure
(\cref{thm:witness-bundle}).  The curvature $F$ is the canonical
obstruction to trivial parallel transport.  The Yang--Mills equation is
the Euler--Lagrange equation of the unique gauge-invariant quadratic
action $\mathrm{tr}(F \wedge *F)$.  Uniqueness of the quadratic action
follows from gauge invariance and dimensional analysis on the $3+1$
base manifold (\cref{thm:spacetime-dimension}).}
\end{proof}
\end{theorem}

\bigskip
\noindent\textbf{Gravity sector (metric modulation).}

\begin{theorem}[Einstein Equation]
\label{thm:einstein}
The witness/Fisher metric $g_{ij}$ on the base manifold $\OmegaCl$ responds
to the distribution of truth-density $\rho$.  The witness action is:
\[
  S[g, \rho] \;=\; \int_{\OmegaCl}
  \Bigl( R(g) + \lambda\, |\nabla \log \rho|^2_g + \mathcal{L}_m \Bigr)\,
  \mathrm{dvol}_g,
\]
where $R(g)$ is the scalar curvature, $\lambda$ is a coupling constant, and
$\mathcal{L}_m$ is the matter Lagrangian.  Variation with respect to $g$
yields the Einstein equation:
\[
  G_{\mu\nu} + \Lambda\, g_{\mu\nu}
  \;=\; T^{(\rho)}_{\mu\nu} + T^{(m)}_{\mu\nu},
\]
where $G_{\mu\nu} = R_{\mu\nu} - \tfrac{1}{2} R\, g_{\mu\nu}$ is the
Einstein tensor, $\Lambda$ is the cosmological constant,
$T^{(\rho)}_{\mu\nu}$ is the truth-density stress-energy, and
$T^{(m)}_{\mu\nu}$ is the matter stress-energy.
\begin{proof}[Proof sketch]
The Fisher metric (Step~30, \cref{thm:riemannian-metric}) on $\OmegaCl$
responds to variations of the truth distribution $\rho$: regions of
high truth-density (many distinguishable truth classes per unit volume)
curve the metric differently from regions of low truth-density.  The
scalar curvature $R(g)$ is the unique scalar invariant of the Riemannian
metric that is second-order in derivatives---forced by diffeomorphism
invariance and dimensional analysis on the $3+1$ base.

The $|\nabla \log \rho|^2_g$ term encodes the energy cost of truth-density
gradients: sharp variations in the density of distinguishable truth classes
cost energy, measured by the Dirichlet form (Step~28).

Variation of $S[g,\rho]$ with respect to $g_{\mu\nu}$ yields the Einstein
equation by the standard Palatini computation.  The Bianchi identity
$\nabla_\mu G^{\mu\nu} = 0$ follows from diffeomorphism invariance of the
witness action---equivalently, from gauge invariance of the truth quotient
(Step~12), which ensures that relabeling spacetime coordinates does not
affect the truth structure.

The cosmological constant $\Lambda$ arises as the vacuum witness energy:
the ground-state energy of the witness algebra $\WitAlg$ acting on $\HilbW$.
Even in the absence of matter, the witness algebra has a nonzero vacuum
expectation value $\langle \OmVec, H\, \OmVec \rangle > 0$ (since
$H = J\nabla E$ is bounded below but not zero on $\OmVec$), contributing a
constant energy density to the right-hand side of the Einstein equation.

\forcing{The Fisher metric is forced by the Dirichlet form (Step~28) and
\v{C}encov's theorem (Step~30).  The Einstein--Hilbert action is the unique
diffeomorphism-invariant scalar action that is second-order in metric
derivatives.  The Bianchi identity is forced by diffeomorphism invariance
(gauge invariance of the truth quotient).  $\Lambda$ is forced as the vacuum
energy of the witness algebra.}
\end{proof}
\end{theorem}

\bigskip
\noindent\textbf{Thermodynamic sector (dissipative).}

\begin{theorem}[Second Law of Thermodynamics]
\label{thm:second-law}
The irreversible part $-\nabla D$ of the split law gives gradient flow
toward equilibrium.  The entropy production rate on the continuum is:
\[
  \frac{dS}{dt} \;=\; \int_{\OmegaCl} |\nabla D|^2\, \mathrm{dvol}_g
  \;\geq\; 0.
\]
\begin{proof}[Proof sketch]
The dissipation functional $D$ on $\OmegaCl$ is non-negative by
construction: it measures the cost of irreversible witness record formation
(Step~8, \cref{thm:time}).  The gradient flow $\dot{x} = -\nabla D(x)$
descends $D$ monotonically, since
$\dot{D} = \langle \nabla D, \dot{x} \rangle = -|\nabla D|^2 \leq 0$.

The entropy $S$ increases because the Riemannian volume swept by the flow
grows: distinguishable truth classes accumulate as records are formed.
The production rate $dS/dt = \int |\nabla D|^2\, \mathrm{dvol}_g \geq 0$
is the continuum extension of Step~8's inequality $\Del\Tim \geq 0$.
Equality $dS/dt = 0$ holds if and only if $\nabla D = 0$
everywhere---equilibrium, where no further irreversible witnessing occurs.

\forcing{The dissipative part of the split law is forced by Step~14.  The
non-negativity of entropy production is forced by the append-only ledger
(Step~5) and irreversibility of time (Step~8).  The integral form extends
the discrete $\Del\Tim \geq 0$ to the continuum $\OmegaCl$.}
\end{proof}
\end{theorem}

% ===========================================================================
\subsection{Empirical Predictions}
\label{sec:physics:predictions}

The physical content derived above yields three falsifiable predictions.
Each identifies a structural consequence of the framework that can be tested
independently of the open computational problems (mass values, coupling
constants).

\begin{description}[leftmargin=2em, labelwidth=2em, labelsep=0.5em]

\item[\textbf{P1.}] \textbf{Inverse-square witness flux.}\;
Any fundamental long-range separative sector must obey inverse-square
witness flux.  Specifically, the conductance/flux matching
(\cref{thm:spacetime-dimension}) requires that every fundamental force
mediating witness separation at long range decays as $r^{-2}$ in the
flux density.  A systematic violation of the inverse-square law at a
truly fundamental scale---not attributable to short-range screening
(as in the weak and strong forces) or finite-size effects---falsifies
the conductance theorem (\cref{lem:conductance-exponent}).

\item[\textbf{P2.}] \textbf{Order-curvature correction for noncommuting
witnesses.}\;
Sequential noncommuting witness instruments must exhibit an
order-curvature correction proportional to the local commutator
$2$-form.  Quantitatively, the probability difference between performing
witness $w_1$ before $w_2$ versus $w_2$ before $w_1$ satisfies:
\[
  \Del P \;\sim\; \omega\bigl(\WitCommutator{w_1}{w_2}\bigr)
  + O\!\bigl(\|\WitCommutator{w_1}{w_2}\|^2\bigr),
\]
where $\omega$ is the symplectic form (Step~32, \cref{thm:symplectic-form})
evaluated on the commutator.  A null result---$\Del P = 0$---where
the ordered witness algebra predicts a nonzero commutator
$\WitCommutator{w_1}{w_2} \neq 0$ falsifies the ordered extension of
A0$^*$ (\cref{def:witness-algebra}).

\item[\textbf{P3.}] \textbf{Dimensionality is a theorem.}\;
The $3+1$ dimensionality result (\cref{thm:spacetime-dimension}) is a
\emph{theorem} of the framework, not an assumption.  Any truly
scale-free long-range witness geometry in $n \neq 3$ spatial dimensions
is incompatible with the conductance/flux compatibility theorem.  This
is falsifiable in the following sense: if a physical system were found
that exhibits fundamental long-range scale-free propagation in an
effectively $n \neq 3$ spatial geometry (not merely a dimensional
reduction or compactification artifact), it would contradict the matching
$n - 1 = 2$ and hence falsify the conductance axioms (SC1)--(SC4).

\end{description}

\begin{remark}[Prediction Scope]
\label{rem:prediction-scope}
Predictions P1--P3 test the \emph{structural} layer of the framework:
dimensionality, flux scaling, and the ordered witness algebra.  The
gauge group $SU(3) \times SU(2) \times U(1)$ is itself derived
(\cref{thm:fiber-decomposition}), not empirical input.
Further predictions---such
as specific mass ratios or coupling constant values---await solution of the
spectral and coarse-graining problems identified in
\cref{rem:mass-boundary,rem:coupling-boundary}.
\end{remark}
